% !PRINTED 656
\weberpage{656}
%<<CONT
\gesperrt{be decomposed into $n$ (equal or different) prime factors, and let}
\weq{16}{\tau_1, \tau_2, \ldots, \tau_n}
\gesperrt{be the constant functionals to which $\tau$, with respect to the
moduli}
\weq{17}{\pi_1, \pi_2, \ldots, \pi_n}
\gesperrt{becomes congruent.}

\gesperrt{Let $\Phi(t)$ be an integral holomorphic functional in $z$
of degree $m$ with respect to $t$, such that the condition}
\weq{18}{
\Phi(\tau) \equiv 0 \quad (\mathrm{mod}\ z - c)
}
\gesperrt{is satisfied. Then:}
\weq{19}{
\Phi(t) \equiv (t - \tau_1)(t - \tau_2)\ldots(t - \tau_n)\,\Psi(t) \quad (\mathrm{mod}\ z - c)\,,
}
\gesperrt{where $\Psi$ is a functional of the same kind of degree $m - n$.}

From this it follows:

5. \gesperrt{A functional $\Phi$ that satisfies the condition (18)
cannot be of lower than degree
$n$.}

And applied to the function $F(t)$, the formula follows:
\weq{20}{
F(t) = N(t - \tau) \equiv (t - \tau_1)(t - \tau_2)\ldots(t - \tau_n) \quad (\mathrm{mod}\ z - c).
}

We can replace the congruence (20) by an equation in the
following way:

As in (17) of the preceding section, we collect in the
decomposition of $z - c$ into prime factors the equal factors
 together, and set:
\weq{21}{
\begin{aligned}
z - c &= \pi_1^{e_1}\pi_2^{e_2}\ldots\pi_m^{e_m}\\
e_1 + e_2 + \cdots + e_m &= n\,,
\end{aligned}
}
then by (20):
\weq{22}{
F(t) = (t - \tau_1)^{e_1}(t - \tau_2)^{e_2}\ldots(t - \tau_m)^{e_m} + (z - c)R(t)\,,
}
where $R(t)$ is an integral functional, in the first place in $\Omega$. But
formula (22) itself says that $R(t)$ is contained in $\overline{Z}$, and
that it is an integral function of the functional variables $t, t_1, t_2,
\ldots, t_n$.

We can therefore also form the derivative of equation (12)
% [sic] printed (12), but (23) is the t-derivative of (22) (exponents e_i from (21), term (z-c)R'(t)), so (22) is meant
with respect to $t$, and obtain:
\weq{23}{
\begin{aligned}
F'(t) = {}& e_1(t - \tau_1)^{e_1-1}(t - \tau_2)^{e_2}\ldots(t - \tau_m)^{e_m}\\
+ {}& e_2(t - \tau_1)^{e_1}(t - \tau_2)^{e_2-1}\ldots(t - \tau_m)^{e_m}\\
& \cdot\ \cdot\ \cdot\ \cdot\ \cdot\ \cdot\ \cdot\ \cdot\ \cdot\ \cdot\ \cdot\ \cdot\ \cdot\ \cdot\\
+ {}& e_m(t - \tau_2)^{e_1}(t - \tau_2)^{e_2}\ldots(t - \tau_m)^{e_m-1} + (z - c)R'(t) % [sic] first factor of last row printed (t-\tau_2) twice, apparently for (t-\tau_1)
\end{aligned}
}
%CONT>>
